\documentclass[manuscript,screen,review,anonymous]{acmart}

\acmJournal{TOCS}

\begin{document}

% ============================================================
% Title
% ============================================================
\title[DynAPM: Drift-Aware Profile Maintenance]{DynAPM: Drift-Aware Profile Maintenance with
                    Predictor-Grounded Forecasting in the
                    Cloud-Edge Continuum}

% ============================================================
% Authors (anonymized for double-anonymous review)
% ============================================================
\author{Anonymous Author(s)}
\renewcommand{\shortauthors}{Anonymous Author(s)}

% ============================================================
% Abstract
% ============================================================
\begin{abstract}
The Cloud-Edge Continuum (CECC) motivates lifecycle management (LCM)
consumers---schedulers, autoscalers, and capacity planners---that read
accurate, current application profiles for placement and provisioning.
Existing profile representations are either static, fixed at deployment,
or refreshed from recent observation windows on a fixed schedule or in
response to threshold crossings; all mis-track applications whose
behavior drifts. We propose DynAPM, a profile-maintenance methodology
intended for CECC LCM systems, whose dynamic fields encode forecasts of
near-future application behavior rather than current state, evaluated
under a snapshot-and-verify protocol with 24-hour ground-truth horizons.
The methodological contribution combines two mechanisms: predictor-grounded
refresh, in which profile fields are populated from per-application
forecasting models rather than recent observations; and drift-augmented
refresh, in which the periodic schedule is augmented with refreshes
triggered by online drift detection on prediction residuals. We instantiate
these mechanisms as Periodic-Predictor and DynAPM-V3 and evaluate them
against four observation-window baselines (Static, Periodic-Window,
Reactive, and Full-Rewrite)---six methods in total---via offline replay of
CPU and memory traces from 1{,}576 production VMs in the Materna
enterprise data-center corpus, representing the data-center/cloud tier of
the continuum. Predictor-grounded refresh improves CPU profile class
accuracy by 13.0 percentage points (95\% confidence interval [12.1, 14.0])
and memory class accuracy by 14.5 percentage points (95\% confidence
interval [13.8, 15.3]) over a matched-schedule observation-window baseline
(paired asymptotic Wilcoxon $p = 5.17 \times 10^{-148}$). On a downstream
resource-provisioning task, predictor-grounded profiles reduce SLA-risk
events by 28\% in relative terms. Drift-augmented refresh adds a small but
statistically robust further improvement (asymptotic Wilcoxon
$p = 2.76 \times 10^{-18}$), with larger gains on high-volatility
workloads. Edge-tier, containerized, heterogeneous, and live-system
validation remain future work.
\end{abstract}

% ============================================================
% CCS Concepts
% ============================================================
\begin{CCSXML}
<ccs2012>
   <concept>
       <concept_id>10010520.10010521.10010537.10003100</concept_id>
       <concept_desc>Computer systems organization~Cloud computing</concept_desc>
       <concept_significance>500</concept_significance>
       </concept>
   <concept>
       <concept_id>10003033.10003068.10003073.10003074</concept_id>
       <concept_desc>Networks~Network resources allocation</concept_desc>
       <concept_significance>300</concept_significance>
       </concept>
   <concept>
       <concept_id>10010147.10010257.10010258.10010259.10010264</concept_id>
       <concept_desc>Computing methodologies~Supervised learning by regression</concept_desc>
       <concept_significance>300</concept_significance>
       </concept>
   <concept>
       <concept_id>10002950.10003648.10003688.10003693</concept_id>
       <concept_desc>Mathematics of computing~Time series analysis</concept_desc>
       <concept_significance>300</concept_significance>
       </concept>
   <concept>
       <concept_id>10010147.10010257.10010282.10010284</concept_id>
       <concept_desc>Computing methodologies~Online learning settings</concept_desc>
       <concept_significance>100</concept_significance>
       </concept>
 </ccs2012>
\end{CCSXML}

\ccsdesc[500]{Computer systems organization~Cloud computing}
\ccsdesc[300]{Networks~Network resources allocation}
\ccsdesc[300]{Computing methodologies~Supervised learning by regression}
\ccsdesc[300]{Mathematics of computing~Time series analysis}
\ccsdesc[100]{Computing methodologies~Online learning settings}

% ============================================================
% Keywords
% ============================================================
\keywords{cloud-edge continuum, application profile, profile maintenance, concept drift, workload forecasting, resource provisioning, lifecycle management, online adaptation}

\maketitle

% ============================================================
% 1. Introduction
% ============================================================
\section{Introduction}\label{sec:intro}

The Cloud-Edge Continuum (CECC) is the operational regime in which
applications may be deployed across a heterogeneous fabric of cloud
regions, edge sites, and on-premises infrastructure, with workloads
flowing dynamically between them in response to demand, latency
requirements, and energy constraints~\citep{ullah_cecc_survey_2023}.
Effective lifecycle management (LCM) in this regime depends on a
continually accurate picture of each application's behavior: how loaded
it is, how variable it is, and what resources it likely needs over the
next operational window. This information is conventionally stored in an
\emph{application profile}---a structured representation of the
application's static identity (provisioning, region, tenant) and its
dynamic behavior (load class, expected resource usage, variability
indicators)~\citep{nfv_sdn_apm}. LCM consumers, including placement
schedulers, autoscalers, capacity planners, and SLA monitors, read these
profiles to make per-application decisions.

The accuracy of these decisions is bounded by the accuracy of the profile.
A profile that says ``this application is currently in a low-load state''
when in reality the next 24 hours will see a sustained burst leads
downstream consumers astray: the autoscaler under-provisions, the placement
engine packs aggressively, the capacity planner over-commits. In production
deployments---including the data-center and cloud tiers that populate
CECC systems---applications drift through behavioral regimes over hours
and days, driven by user activity, batch schedules, and external events.
A profile representation that does not adapt to this drift becomes
progressively misaligned with reality.

Existing approaches to profile maintenance are limited in two ways.
\emph{Static profiles}, set at application deployment and never updated,
miss all post-deployment drift. \emph{Reactive update policies}, which
refresh profile fields when an instantaneous metric crosses a threshold,
chatter under noise and miss slow drift. \emph{Periodic refresh policies},
which update profile fields on a fixed schedule (e.g., every six hours)
from a recent observation window, are more stable but ground their update
content in past observations only---they extrapolate from what just
happened to predict what comes next, even when a per-application
forecasting model could provide a more informed prediction. The space of
profile-maintenance designs is large, but the question of \emph{where}
the update content should come from---whether observed history or
model-based forecast---has received little systematic attention in the
LCM literature.

This paper addresses that gap. We propose DynAPM, a profile-maintenance
methodology intended for CECC LCM systems, in which dynamic fields encode
forecasts of near-future application behavior, validated under a
snapshot-and-verify protocol with 24-hour ground-truth horizons. The
framework combines two refresh mechanisms: predictor-grounded periodic
refresh, in which scheduled updates are populated from per-application
forecasting models rather than recent observation windows; and
drift-augmented refresh, in which the periodic schedule is supplemented
by refreshes triggered by an online drift detector running on the
prediction residual stream. DynAPM is not a predictive autoscaler: it
maintains a persistent, multi-consumer profile artifact rather than
directly issuing scaling actuations.

We evaluate DynAPM via offline replay of CPU and memory traces from
1{,}576 production VMs in the Materna enterprise data-center corpus,
which we treat as representative of the data-center/cloud tier of the
continuum. We compare two proposed methods---Periodic-Predictor and
DynAPM-V3---against four observation-window baselines (Static,
Periodic-Window, Reactive, and Full-Rewrite): six methods in total. The
evaluation uses an explicit snapshot-and-verify protocol with paired
bootstrap confidence intervals and asymptotic Wilcoxon signed-rank tests.
To assess operational consequence, we additionally evaluate a downstream
resource-provisioning recommendation task in which an LCM consumer
recommends provisioned capacity from each method's profile snapshot and
we measure the resulting rate of SLA-risk events. We do not claim
empirical coverage of edge-tier, containerized, or live orchestrated
deployments; those settings remain future work
(Section~\ref{sec:discussion-threats}).

The contributions of this paper are:
\begin{itemize}
\item \textbf{Predictor-grounded profile refresh:} A methodology in which
profile dynamic fields are populated from per-application forecasting
models rather than recent observation windows. On 1{,}576 VMs, this
single change improves CPU profile class accuracy by 13.0 percentage
points (95\% CI [12.1, 14.0]) and memory class accuracy by 14.5
percentage points (95\% CI [13.8, 15.3]) over a matched-schedule
observation-window baseline (paired asymptotic Wilcoxon
$p = 5.17 \times 10^{-148}$).

\item \textbf{Drift-augmented refresh:} An online refinement in which
scheduled refreshes are augmented by drift-triggered refreshes, both
grounded in the predictor's forecast. Drift augmentation adds a small
but statistically robust further improvement of 0.24 percentage points
CPU and 0.32 percentage points memory class accuracy (asymptotic
Wilcoxon $p = 2.76 \times 10^{-18}$), with effect concentrated on
high-volatility workloads and no observed degradation on any tier.

\item \textbf{Snapshot-and-verify evaluation protocol:} A methodology
for evaluating profile-shaped representations in which a profile snapshot
at time $t$ is verified against ground truth derived from $[t, t+24\mathrm{h})$
once that horizon has streamed past. We provide paired bootstrap
confidence intervals and asymptotic Wilcoxon signed-rank tests for all
headline comparisons.

\item \textbf{Operational impact via downstream provisioning:} A direct
evaluation of how profile accuracy translates to operational outcomes.
Predictor-grounded profiles reduce SLA-risk events (under-provisioning)
by 28\% in relative terms and 2.4 percentage points in absolute terms.
\end{itemize}

These contributions are methodological and are empirically supported on
production data-center VMs; they are intended for CECC LCM consumers but
are not validated here across edge, container, or live multi-tier systems.



The remainder of this paper is organized as follows.
Section~\ref{sec:bg} surveys related work in cloud workload forecasting,
reactive autoscaling, concept drift adaptation, and application profiling
for the CECC. Section~\ref{sec:framework} formalizes the DynAPM
framework. Section~\ref{sec:setup} describes the experimental setup.
Section~\ref{sec:results} presents the corpus-scale results.
Section~\ref{sec:discussion} discusses implications and threats to
validity. Section~\ref{sec:conclusion} concludes.

% ============================================================
% 2. Background and Related Work
% ============================================================
\section{Background and Related Work}\label{sec:bg}

The problem of maintaining accurate, up-to-date application
profiles in dynamic compute environments draws on four
established research areas: cloud workload forecasting, reactive
resource adaptation, concept-drift detection in streaming data,
and application profiling for the cloud-edge continuum. Each of
these literatures has matured along its own trajectory, and
each contributes something to the design of an adaptive
profile-maintenance method, but no prior work has integrated
them around the central question we pose: \emph{where} should
the content of profile updates come from? We review each area
in turn, drawing out the threads that converge on this question
and clarifying how our contribution sits relative to existing
work.

Workload forecasting in cloud and data-center environments is
the most mature of the four. Classical statistical approaches
such as ARIMA were among the earliest methods applied to
cloud-resource prediction, with Calheiros et al.\ demonstrating
that ARIMA-based prediction can achieve up to 91\% accuracy on
web-server workloads and translate into measurable
quality-of-service gains~\cite{calheiros_arima_2015}. With the
availability of large-scale public traces such as the Google
Borg traces~\cite{tirmazi_borg_2020} and the Alibaba production
traces~\cite{lu_alibaba_2017}, the community has shifted toward
machine-learning approaches that exploit per-VM and
per-application heterogeneity. More recent work emphasizes deep
models and hybrid architectures: LSTM and BiLSTM networks have
been extensively applied to workload and energy-state
forecasting at the VM
level~\cite{khan_workload_forecasting_2022}, and hybrid models
that combine signal decomposition, attention mechanisms, and
LSTM variants achieve further accuracy gains on Google and
Alibaba traces~\cite{bi_hybrid_lstm_2023}. RandomForest
predictors~\cite{breiman_random_forests_2001} remain a
competitive non-deep alternative, particularly when moderate
training cost on tabular features is preferred over deep sequence
models. Across this literature, the
dominant evaluation metrics are per-step prediction quality
measures: $R^2$, mean absolute error, and root mean squared
error, reported on held-out test windows. This methodology
answers the question of how accurate a model's prediction of
the next observation is, but does not address how well that
prediction serves a downstream operational consumer. DynAPM
uses a RandomForest forecaster as a component, but our
methodological focus differs: we evaluate the maintained
profile representation that a lifecycle-management consumer
would actually read, rather than the per-step accuracy of the
underlying forecaster.

Workload forecasters do not exist in isolation; they feed into
the resource-adaptation systems that act on their predictions.
A separate strand of work has therefore addressed online
adaptation of resource allocations in response to observed
workload changes. Reactive autoscalers monitor instantaneous
metrics (CPU utilization, request queue depth, response
latency) and trigger scale-up or scale-down events when
thresholds are crossed. Production systems such as Kubernetes
Horizontal Pod Autoscaler~\cite{kubernetes_hpa} and AWS Auto
Scaling~\cite{aws_autoscaling} implement variants of this
pattern. Threshold-based reactive policies are simple and
well-understood, but they suffer from two known problems:
oscillation under noise, often described as chatter, and lag
in response to sustained changes. Periodic update policies
represent a complementary design point: rather than triggering
on individual observations, they refresh state at fixed
intervals using aggregate statistics over a recent window.
Both reactive and periodic strands share an implicit
assumption that the input to any update decision is the
recently observed resource trace.

A substantial body of work goes further and uses forecasts
to \emph{proactively} change resource allocations---predictive
or proactive autoscaling. Surveys of cloud auto-scaling
distinguish reactive threshold policies from proactive
approaches based on time-series prediction, control theory,
and related methods~\cite{lorido_botran_autoscaling_2014,qu_autoscaling_survey_2018}.
Concrete systems couple demand predictors to scaling
actuators: PRESS predicts resource demand and adjusts
allocations to reduce waste and SLO risk~\cite{gong_press_2010};
prediction models have likewise been used for adaptive cloud
resource provisioning ahead of demand~\cite{islam_prediction_provisioning_2012}.
At production scale, Google Autopilot continuously adjusts
horizontal replica counts and vertical CPU/memory limits from
historical usage~\cite{rzadca_autopilot_2020}, and hybrid
frameworks such as OptScaler combine proactive forecasts with
reactive feedback under model-predictive control to decide
scaling actions~\cite{zou_optscaler_2024}. Across this
literature, the system output is an \emph{actuation}---a replica
count, a capacity limit, or a provisioning decision---rather than
a durable, multi-consumer description of application behavior.
DynAPM's predictor-grounded refresh is complementary: at each
scheduled refresh, we ground a shared profile artifact in the
predictor's forecast of the next 24 hours rather than in the most
recent 24 hours of observations, without directly issuing scaling
actions. The architectural change is small but, as
Section~\ref{sec:results} shows, empirically consequential.

A complementary line of research addresses the case in which
the workload itself changes its statistical character over
time --- the problem of concept drift in streaming data. The
canonical methods are mature: sequential change-point tests
(CUSUM)~\cite{page_cusum_1954}, the Page-Hinkley
test~\cite{hinkley_changepoint_1971}, and ADWIN (Adaptive
Windowing)~\cite{bifet_adwin_2007}. ADWIN maintains a
variable-length window over a numerical stream and signals
drift when sub-windows exhibit statistically distinct means,
with formal guarantees on its false-positive rate that have
made it a common default in streaming-ML toolkits such as
River~\cite{montiel_river_2021}. The application of drift
detection specifically to cloud workloads is more recent.
Adaptive sliding-window techniques have been applied to
data-center resource utilization estimation, demonstrating
that window adaptation based on observed change rate improves
utilization prediction over fixed-window
baselines~\cite{baig_adaptive_windows_2019}. We use ADWIN in
a different role from these prior works: not as a signal that
invalidates a model and triggers retraining, but as one of two
refresh triggers for a forecast-grounded profile. Drift events
in our design augment a fixed periodic schedule with
additional refreshes when the residual stream signals change
between scheduled refreshes; the underlying forecaster remains
frozen.

These three literatures --- forecasting, reactive adaptation,
and drift detection --- all concern the dynamics of resource
management. The fourth concerns the data structure that
mediates between resource managers and application instances:
the application profile itself. The cloud-edge continuum
extends classical cloud computing to a distributed fabric in
which compute may execute at central clouds, regional edge
sites, or on-premises infrastructure, and recent surveys
characterize the orchestration challenges this distributed
fabric introduces, including resource heterogeneity, mobility
of data sources, and lifecycle complexity for applications
spanning multiple tiers~\cite{ullah_cecc_survey_2023}.
Lifecycle management in this regime requires structured
representations of applications and their requirements that
downstream consumers --- placement schedulers, autoscalers,
SLA monitors --- can query at decision time. The
structured-profile pattern itself appears in network function
virtualization, where Application Profile Management schemas
formalize what static and dynamic fields a profile carries and
how consumers retrieve them. A unified APM template for
microservices-based applications in the cloud-edge continuum
has previously been proposed~\cite{nfv_sdn_apm}, defining
static identity fields and dynamic behavior fields with formal
OWL-based semantics. The present paper builds on that schema
by addressing a complementary question: not what fields a
profile should carry, but how those fields should be kept up
to date as the underlying workload changes.

The contribution of this paper sits at the intersection of
these areas but is distinct from each. Relative to
workload forecasting, we do not propose a novel forecasting
model; we use a standard RandomForest predictor and ask a
different question---what is the right way to use a
forecaster's output as the source of profile updates. Relative
to reactive and predictive autoscaling, we evaluate the
profile-maintenance layer that sits beneath such consumers and
characterize the accuracy implications of different update
sources at that layer, rather than issuing scaling actuations.
Relative to drift-detection literature, we use ADWIN in
a refresh-augmentation role rather than as a model-invalidation
signal, and we measure whether the augmented refresh improves
an operationally meaningful metric. Relative to CECC profiling
schemas, we treat the continuum as the motivating deployment
context for LCM consumers of application profiles, but we do
not claim a multi-tier continuum evaluation. The remaining gap
is a controlled empirical comparison of \emph{profile update
sources} for adaptive lifecycle management. DynAPM addresses
that gap with forecast-based profile semantics, a
snapshot-and-verify protocol, and a corpus-scale offline
evaluation on 1{,}576 production data-center VMs.

% ============================================================
% 3. The DynAPM Framework
% ============================================================
\section{The DynAPM Framework}\label{sec:framework}

This section formalizes the DynAPM framework: the profile schema
and what its dynamic fields encode (\ref{sec:framework-schema}),
the snapshot-and-verify protocol used to evaluate any
profile-maintenance method (\ref{sec:framework-verify}), the
per-application forecasting model (\ref{sec:framework-forecaster}),
and the two proposed profile-maintenance methods:
Periodic-Predictor (\ref{sec:framework-pp}) and DynAPM-V3 with
drift augmentation (\ref{sec:framework-v3}). Four observation-window
baselines are described in Section~\ref{sec:framework-baselines}.

\subsection{Profile Schema and Forecast Semantics}\label{sec:framework-schema}

A DynAPM profile is a key--value record associated with a single
application instance (here, a single virtual machine). The schema
extends the unified APM template proposed in
prior work~\cite{nfv_sdn_apm} by adding forecast-grounded dynamic
fields. The schema partitions fields into two groups.

\textbf{Static fields} are fixed at deployment and unchanged
during the application's lifecycle: an application identifier,
provisioned CPU and memory capacity, deployment region, tenant,
and similar metadata. These fields are written once and read
many times; their maintenance is outside the scope of this work.

\textbf{Dynamic fields} are intended to summarize the application's
near-future operational behavior. In the present formulation,
the dynamic fields are:
\begin{itemize}
\item \texttt{cpu\_class} and \texttt{memory\_class}: discrete load
classes drawn from the five-element label space \{\textsc{idle},
\textsc{low}, \textsc{medium}, \textsc{high}, \textsc{burst}\},
defined by per-application empirical quantile thresholds.
\item \texttt{predicted\_cpu\_p95\_next\_24h} and
\texttt{predicted\_mem\_p95\_next\_24h}: numerical 95th-percentile
forecasts of CPU and memory utilization over the next 24-hour window.
\end{itemize}

The class fields and the p95 fields are independent views of the
same forecast: the numerical p95 is computed directly as the 95th
percentile of the predicted next-24h utilization window, while the
class label is computed by mapping the mean of that same predicted
window through the per-application class thresholds. This decoupling
is deliberate. Downstream LCM consumers concerned with absolute
provisioning capacity read the p95 fields directly; consumers that
operate over discrete behavioral categories (e.g., placement
schedulers that reason in terms of load tiers) read the class
fields. Both fields are written together at every refresh.

The semantics of these fields differ from prior practice in a key
respect. In conventional periodic update policies, the dynamic
fields summarize \emph{recently observed} behavior --- for example,
the empirical p95 over the previous 24 hours and the corresponding
class derived from the same window. In DynAPM, the dynamic fields
summarize \emph{predicted future} behavior over a forward 24-hour
horizon. The class label and the numerical p95 are both derived
from the forecaster's output rather than from observed history.

This forecast-grounded semantics makes the profile directly
actionable for downstream consumers. A capacity planner that
provisions for the next 24 hours can use
\texttt{predicted\_cpu\_p95\_next\_24h} directly as input to a
provisioning decision (e.g., $\mathrm{capacity} =
\texttt{predicted\_cpu\_p95\_next\_24h} \times s$, where $s$ is
a safety factor). Under window-based semantics, the consumer
would instead receive an estimate of past behavior and would have
to perform its own forecasting, fragmenting forecast logic across
multiple LCM components.

The class thresholds for an application $A$ are computed from the
bootstrap segment (Section~\ref{sec:framework-forecaster}) as the
empirical 20th, 40th, 60th, and 80th percentiles of $A$'s
utilization on that segment. Computing these per-application
yields a stable per-VM calibration of the five-class label space:
$T_{0.20}^A, T_{0.40}^A, T_{0.60}^A, T_{0.80}^A$ partition the
real line into five intervals corresponding to \textsc{idle},
\textsc{low}, \textsc{medium}, \textsc{high}, and \textsc{burst}.
CPU and memory thresholds are computed independently from the same
bootstrap segment but using their respective metric streams; they
typically differ in absolute value.

\subsection{Snapshot-and-Verify Evaluation Protocol}\label{sec:framework-verify}

To evaluate any profile-maintenance method, we record what the
profile says at a sequence of evaluation timestamps and compare
those snapshots against ground truth derived from the streaming
data that follows.

Let $A$ denote an application with utilization stream
$\{u_t\}_{t \geq 0}$. A profile-maintenance method $M$ produces,
at each timestamp $t$, a profile snapshot $\Pi^M_A(t)$ containing
the four dynamic fields described above. Evaluation proceeds at a
fixed grid of evaluation timestamps drawn from the streaming
segment of the trace; at each evaluation timestamp $t_e$, the
method's most recently committed snapshot $\Pi^M_A(t_e)$ is
recorded. The exact cadence of this evaluation grid and its
data-aware filtering criterion are specified in
Section~\ref{sec:setup-eval}; importantly, the evaluation cadence
is independent of the refresh cadence used by any individual
method, so a method that refreshes more often than evaluation
occurs is read at evaluation time with whatever snapshot it has
most recently written.

\textbf{Ground truth} is derived from the future window
$[t_e, t_e + 24\mathrm{h})$ once that window has streamed past. The
ground-truth class is computed by mapping the mean of observed
utilization over that same forward window
$[t_e, t_e + 24\mathrm{h})$ through the same per-application class
thresholds $T_q$ used for the profile's own class derivation. The
ground-truth p95 is the empirical 95th percentile of utilization
over $[t_e, t_e + 24\mathrm{h})$. Class and p95 ground truth
therefore share the same forward horizon as the predicted fields.

\textbf{Per-snapshot accuracy} is then computed as
$\mathbb{1}[\texttt{cpu\_class}(\Pi^M_A(t_e)) =
\mathrm{class}_\mathrm{gt}(t_e)]$ for the class metric, and
$|\texttt{predicted\_cpu\_p95}(\Pi^M_A(t_e)) -
\mathrm{p95}_\mathrm{gt}(t_e)|$ for the p95 absolute-error metric.
Per-VM accuracies are averaged over evaluation timestamps; corpus
accuracies are averaged over VMs. Pairwise statistical
comparisons (Section~\ref{sec:setup-statistics}) operate at the
per-VM level.

A consequence of this protocol is that evaluation requires a
warm-up period at both ends of the trace: the first 30\% of each
VM's time series is reserved as the bootstrap segment for the
per-application forecasting model and class thresholds (see
Section~\ref{sec:framework-forecaster}), and the last 24 hours of
the streamed period cannot be verified because the forward window
extends past the trace.

\subsection{Per-Application Forecasting Model}\label{sec:framework-forecaster}

DynAPM uses a per-application forecasting model trained offline
and frozen during operational evaluation. We use a RandomForest
regressor~\cite{breiman_random_forests_2001} as a practical
per-VM baseline on the engineered tabular feature set of lags,
rolling-window statistics, and calendar encodings described below.
This choice fits the study's constraints: moderate training cost
on CPU-only hardware, low inference complexity for repeated
lookahead generation, and suitability for fitting separate models
on each VM's bootstrap segment without specialized accelerators.
We do not claim that RandomForest is uniquely robust, efficient,
or hyperparameter-insensitive; hyperparameters are fixed a priori
and are not tuned in this work.

For each application $A$, we reserve a per-application
\emph{bootstrap segment} consisting of the first 30\% of the
streaming time series. The 30\% fraction is fixed across all VMs;
the corresponding duration in days varies per VM according to the
length of its time series. Two regressors are trained on this
segment, one for CPU usage and one for memory usage, each with
the same input features. The features are built exclusively from
the target metric's history at the same VM and include lagged
values (at lags 1, 2, 3, 6, 12, 24, and 48 steps),
rolling-window statistics (mean, standard deviation, minimum,
maximum, and range over windows of 6, 12, 24, 48, and 96 steps),
first and seasonal differences (lags 1 and 12), and cyclical
hour-of-day and day-of-week encodings derived from the timestamp
index. The first 96 timesteps of feature engineering are dropped
before training, since the longest rolling-window features are
not defined on earlier timestamps. The target variable is the
next-step value of the same metric; the regressor is therefore a
one-step-ahead numerical forecaster.

Once trained, the regressor is frozen for the remainder of the
trace. To produce the 24-hour forecasts that the profile-maintenance
methods consume, the frozen regressor is rolled forward 288 steps
(corresponding to 24 hours at the Materna trace's 5-minute sampling
cadence) starting from the current feature vector. The resulting
288-element predicted window provides both the numerical p95
(via the 95th percentile of the window) and the input to the class
derivation (via the window mean mapped through the per-application
class thresholds).

We deliberately do not propose a novel forecasting architecture.
The contribution of DynAPM is in how forecaster output is consumed
by the profile-maintenance layer, not in the forecaster itself.
Replacing RandomForest with a more expressive model is
straightforward and is left as future work.

\subsection{Method: Periodic-Predictor}\label{sec:framework-pp}

The Periodic-Predictor method refreshes the dynamic fields of the
profile on a fixed periodic schedule, with each refresh's content
sourced from the per-application forecaster's current output.

Concretely, at every $\Delta_p$-hour interval (we use $\Delta_p = 6$
hours throughout this work), the method invokes the forecaster on
the current feature vector, rolls it forward 288 steps to obtain
a predicted 24-hour window, computes the p95 (numerical field) and
the window mean (mapped through per-application class thresholds
to produce the class field), and writes all four dynamic fields to
the profile in a single update event.

Algorithmically:
\begin{verbatim}
on every periodic tick t_p:
    fwin_cpu = roll_forward(forecaster_cpu, features(t_p), n=288)
    fwin_mem = roll_forward(forecaster_mem, features(t_p), n=288)
    Pi.predicted_cpu_p95_next_24h = p95(fwin_cpu)
    Pi.predicted_mem_p95_next_24h = p95(fwin_mem)
    Pi.cpu_class    = class_of(mean(fwin_cpu), T_cpu)
    Pi.memory_class = class_of(mean(fwin_mem), T_mem)
    log_refresh(t_p, source="periodic")
\end{verbatim}

Between periodic ticks, the profile is unchanged. The schedule is
deterministic and independent of the observed utilization stream.

The contrast with conventional window-based periodic refresh is
narrow but consequential. A window-based method, given the same
schedule, would compute the dynamic fields from the previous 24
hours of \emph{observed} utilization. Periodic-Predictor instead
sources its dynamic fields from the predictor's forecast of the
\emph{next} 24 hours. The schedule is identical; only the source
of the update content differs. Section~\ref{sec:results} shows
that this single change improves per-snapshot class accuracy by
13.0 percentage points for CPU and 14.5 percentage points for
memory on 1{,}576 VMs.

\subsection{Method: DynAPM-V3 with Drift Augmentation}\label{sec:framework-v3}

DynAPM-V3 augments Periodic-Predictor with online drift detection
that can trigger additional refreshes between scheduled periodic
ticks. The method maintains \emph{two} independent ADWIN
detectors~\cite{bifet_adwin_2007}, one for CPU and one for memory,
implemented via River~\cite{montiel_river_2021} with significance
parameter $\delta = 0.002$ and River's remaining defaults
(notably $\mathrm{clock}=32$).

At each streaming timestamp $t$, residuals are formed from the
frozen per-metric RandomForest streaming predictions
(precomputed offline), not from the profile's stored dynamic
fields:
\begin{align}
r^{\mathrm{cpu}}_t &= u^{\mathrm{cpu}}_t - \hat{u}^{\mathrm{cpu}}_t, \\
r^{\mathrm{mem}}_t &= u^{\mathrm{mem}}_t - \hat{u}^{\mathrm{mem}}_t,
\end{align}
where $u^{\cdot}_t$ is observed utilization and
$\hat{u}^{\cdot}_t$ is the corresponding RF prediction at $t$.
Each residual is fed to its own ADWIN instance. A drift event is
declared if \emph{either} detector signals change. That event
triggers a \emph{joint} profile refresh: all four dynamic fields
are rewritten from the CPU and memory predictor lookaheads
exactly as in Periodic-Predictor (next up to 288 steps of
$\hat{u}$ starting at $t$). If a periodic tick is also due at
$t$, the periodic refresh takes precedence and is logged as
\texttt{periodic}; otherwise the refresh is logged as
\texttt{drift}. Drift refreshes do not advance the periodic
timer, so the 6-hour schedule remains calendar-aligned.

DynAPM does not call an explicit detector reset after a refresh.
River ADWIN fully resets a detector on the subsequent
\texttt{update} after it has reported \texttt{drift\_detected};
CPU and memory detectors therefore recover independently
according to which stream signaled. Profile fields persist until
the next refresh (no forecast TTL). Near the end of a VM series,
the lookahead may contain fewer than 288 steps; an empty
lookahead yields no update.

Algorithmically:
\begin{verbatim}
initialize:
    adwin_cpu <- ADWIN(delta=0.002)   # River defaults otherwise
    adwin_mem <- ADWIN(delta=0.002)
    last_periodic <- None             # first streaming step is due

on every streaming timestamp t:
    r_cpu <- u_cpu[t] - yhat_cpu[t]   # RF streaming pred, not profile
    r_mem <- u_mem[t] - yhat_mem[t]
    cpu_drift <- adwin_cpu.update(r_cpu)
    mem_drift <- adwin_mem.update(r_mem)

    periodic_due <- (last_periodic is None)
                    or (t - last_periodic >= Delta_p)   # Delta_p = 6h
    drift_event  <- cpu_drift or mem_drift

    if periodic_due:
        refresh_all_fields_from_predictor_lookahead(t, n<=288)
        log_refresh(t, source="periodic")
        last_periodic <- t
    else if drift_event:
        refresh_all_fields_from_predictor_lookahead(t, n<=288)
        log_refresh(t, source="drift")
        # do not advance last_periodic
\end{verbatim}

The two refresh sources---periodic and drift---both write the
same field set, both use the same forecaster, and both produce
fully-formed profile updates. They differ only in their trigger.
Periodic refreshes occur at fixed cadence; drift refreshes occur
when either residual stream signals change between scheduled ticks.

The mechanism is intentionally conservative on update frequency.
On the corpus reported in Section~\ref{sec:results}, drift events
add approximately 6 additional refreshes per VM beyond the
approximately 85 periodic refreshes per VM---a 7\% increase in
refresh activity. The improvement in profile accuracy that this
additional activity produces is small in absolute magnitude
(0.24 percentage points CPU class accuracy) but statistically
robust (asymptotic Wilcoxon $p = 2.76 \times 10^{-18}$) and
concentrated on high-volatility workloads where drift events are
most informative.

\subsection{Baseline Methods}\label{sec:framework-baselines}

We compare DynAPM-V3 and Periodic-Predictor against four baseline
profile-maintenance methods that span the design space.

\paragraph{Static.} Profile dynamic fields are populated once at the
end of the bootstrap segment, using the empirical p95 and class of
the bootstrap window, and are never updated thereafter. This
represents the worst case for drift handling but is the design
point that current LCM systems often default to in practice when
no maintenance policy is specified.

\paragraph{Periodic-Window.} Profile dynamic fields are refreshed on
the same $\Delta_p = 6$-hour schedule as Periodic-Predictor, but
each refresh's content is sourced from the previous 24-hour
\emph{observation} window rather than from the forecaster's
prediction. The numerical p95 field is the empirical 95th
percentile of that observation window; the class field is derived
by mapping the window's mean through the per-application class
thresholds. This is the most direct comparison point for
Periodic-Predictor: identical schedule, identical update event
structure, identical class-and-p95 derivation logic; only the
source of the underlying utilization window differs (past
observation versus predicted future).

\paragraph{Reactive.} Profile dynamic fields are refreshed when the
class derived from a rolling 6-hour observation window changes
relative to the class currently committed in the profile.
Specifically, at each streaming timestep, the method computes the
mean of the most recent 6 hours of observed CPU (and separately,
memory) utilization, maps each mean through the per-application
class thresholds, and triggers a profile refresh if either the
candidate CPU class or the candidate memory class differs from the
currently committed class. When triggered, the method commits both
modalities together, sourcing the new dynamic fields from the
observation window. This pattern represents the threshold-based
reactive autoscaling design point exemplified by production
systems such as Kubernetes HPA~\cite{kubernetes_hpa} and AWS Auto
Scaling~\cite{aws_autoscaling}, mapped to the
profile-maintenance layer; the design is purely observation-based
and does not consult the forecaster.

\paragraph{Full-Rewrite.} Profile dynamic fields are recomputed on
every observation timestep from the trailing 24-hour
\emph{observation} windows for CPU and memory (the same
window-aggregation helpers used by Periodic-Window and Reactive).
Predictor streams are accepted by the evaluation runner but are
ignored by this baseline. It represents an observation-grounded
upper bound on update frequency: a profile that is always
maximally fresh with respect to recent observations, at the cost
of substantial write churn. This baseline is included to test
whether observation-window methods can close the accuracy gap to
predictor-grounded refresh merely by refreshing more often.

The six methods together span the trade-off space along two
dimensions: source of update content (observation window versus
forecaster output) and refresh trigger (none, threshold-based,
fixed schedule, every-step, or drift-augmented). The four
baselines are observation-window methods. Periodic-Predictor and
DynAPM-V3 occupy a previously unexamined region: forecaster-sourced
content combined with stable refresh policies (fixed schedule,
optionally drift-augmented).

% ============================================================
% 4. Experimental Setup
% ============================================================
\section{Experimental Setup}\label{sec:setup}

This section describes the dataset, forecasting-model
configuration, evaluation protocol parameters, statistical
methodology, and computing environment used to evaluate DynAPM.

\subsection{Dataset}\label{sec:setup-dataset}

We use the Materna VM utilization trace, distributed via the
Grid Workloads Archive~\cite{iosup_gwa_2008,materna_trace}. The
trace consists of three monthly snapshots from a distributed
data center operated by Materna GmbH (Dortmund, Germany), each
covering a one-month window over a heterogeneous set of business
applications running on a VMware ESX
environment~\cite{kohne_federatedcloudsim_2014,kohne_sla_vm_2016}.
Per-VM time series are sampled every 5 minutes and report 11
fields including provisioned CPU and memory, observed CPU and
memory usage (in MHz, percentage, and KB units), disk I/O
throughput, and network throughput. We use only the CPU and
memory percentage fields in the evaluations reported here.

The three traces collectively contain 1{,}592 distinct VM time
series. We retain only VMs whose time series contain at least
2{,}000 observations (corresponding to approximately 7 days of
5-minute samples). This filter removes 16 short series, yielding
a final corpus of 1{,}576 VMs: 511 from Materna-Trace-1, 521 from
Materna-Trace-2, and 544 from Materna-Trace-3.

Materna therefore supplies a production data-center/cloud-tier
workload under offline trace replay. It does not include edge-tier
devices, container or Kubernetes workloads, heterogeneous
cloud--edge testbeds, or live orchestrator integration; those
settings are outside the empirical scope of this paper.

For each retained VM, we partition the time series temporally:
the first 30\% of timestamps form the \emph{bootstrap segment},
used to compute per-application class thresholds and to train
the forecasting model; the remaining 70\% form the \emph{streaming
segment}, used for evaluation. The 30\% fraction is fixed across
all VMs; the corresponding bootstrap duration in days varies per
VM with a mean of 9.3 days, median of 9.3 days, and a range of
2.1 to 15.6 days (standard deviation 1.5 days).

\subsection{Volatility Stratification}\label{sec:setup-volatility}

To support per-tier analysis, we stratify the corpus by workload
volatility. For each VM, we compute the population standard
deviation of CPU utilization (in percentage points) over its
bootstrap segment. This per-VM scalar is a stable, easily-computed
proxy for the magnitude of behavioral variation in the workload.
We then partition the corpus into three tiers --- \textsc{low},
\textsc{medium}, and \textsc{high} --- using the empirical
tertiles of this distribution as cut points. Tier sizes are
525, 525, and 526 VMs respectively. The tier assignment is fixed
at corpus-build time and does not change during evaluation.

\subsection{Per-Application Forecasting Model}\label{sec:setup-forecaster}

The per-VM forecasting model is a RandomForest
regressor~\cite{breiman_random_forests_2001} from scikit-learn,
configured with 200 trees, maximum depth 20, square-root feature
subsampling at each split, and a fixed random seed (42). All
other RandomForest hyperparameters take scikit-learn defaults.
These settings were fixed before corpus evaluation; we do not
report a hyperparameter search or sensitivity study. Each VM
receives two independently trained regressors: one for CPU usage
percentage and one for memory usage percentage. Training uses the
full bootstrap segment after the feature construction in
Section~\ref{sec:framework-forecaster}. No internal validation
split is performed within the bootstrap segment: the forecaster
is a frozen component for the profile-maintenance comparison, not
an object of model selection in this study. The bootstrap segment
contains roughly 2{,}600 samples on average after the corpus
length filter.

The first 96 timesteps of feature engineering are dropped before
training, because the longest rolling-window features (96-step
rolling statistics) are not defined on earlier timestamps. After
this trimming, rows containing any NaN are also dropped. Once
trained, the regressor is frozen for the remainder of the trace.
At each refresh during the streaming segment, the method consumes
the precomputed predictor lookahead of up to 288 steps to form the
24-hour predicted window from which the p95 (numerical field) and
the mean (class derivation) are computed.

\subsection{Evaluation Protocol Parameters}\label{sec:setup-eval}

The snapshot-and-verify protocol introduced in
Section~\ref{sec:framework-verify} is parameterized as follows.

Evaluation timestamps are drawn from a fixed 12-hour grid
beginning at the start of each VM's streaming segment, with the
constraint that the forward 24-hour window from each grid point
must contain at least 200 observations (corresponding to roughly
70\% sampling density at 5-minute cadence). This data-aware
filter removes grid points where the trace is too sparse to
support reliable ground-truth computation. Across the 1{,}576-VM
corpus, the protocol generates 40.5 verified evaluation points
per VM on average (median 41, range 7--48, standard deviation
7.0).

Ground truth at each evaluation timestamp $t_e$ is derived as
follows: both the ground-truth class and the ground-truth p95 are
computed from the forward 24-hour window
$[t_e, t_e + 24\mathrm{h})$. The class is the mean of observed
utilization over that window mapped through the per-application
class thresholds; the p95 is the empirical 95th percentile over
the same window. Snapshot accuracy at each evaluation
timestamp is computed as exact-match for the class metric and
absolute error in percentage points for the p95 metric.

The evaluation cadence (12 hours) is independent of the refresh
cadence used by Periodic-Predictor and DynAPM-V3 (6 hours), and
of the refresh policies of the other baselines. At each
evaluation timestamp, the method's most recently committed
profile snapshot is read; consequently every method is evaluated
on the same set of timestamps for the same VM, eliminating
cadence-mismatch as a source of comparison bias.

\subsection{Provisioning Recommendation Task}\label{sec:setup-provisioning}

The downstream provisioning task uses the same evaluation
timestamps. At each timestamp, the LCM consumer reads the
current profile snapshot and computes a recommended capacity as
$c = \texttt{predicted\_cpu\_p95\_next\_24h} \times s$, where $s$
is a safety factor. We set $s = 1.2$ as the primary configuration
and report sensitivity at $s \in \{1.1, 1.5\}$.

The recommendation is compared against the actual observed p95
of the same forward 24-hour window. We classify the outcome at
each evaluation point into three categories:
\textbf{under-provisioned} ($\mathrm{actual} > c$), corresponding
to an SLA-risk event; \textbf{over-provisioned}
($c > \mathrm{actual} \times 1.5$), corresponding to a wasted
capacity event; and \textbf{well-provisioned} otherwise. We
report per-method rates of each outcome category, averaged first
within each VM across its evaluation points, then across VMs.

\subsection{Statistical Methodology}\label{sec:setup-statistics}

For each pairwise method comparison, we compute three statistics
on the per-VM means.

\paragraph{Mean difference.} The mean of paired per-VM differences
between methods is reported with a 95\% confidence interval
computed via 1{,}000-sample paired bootstrap with replacement
over the VM index, taking the 2.5th and 97.5th percentiles of
the bootstrapped means.

\paragraph{Wilcoxon signed-rank test.} The two-sided paired Wilcoxon
signed-rank test on per-VM differences, implemented with SciPy's
\texttt{wilcoxon} using \texttt{zero\_method="wilcox"} (zero
differences discarded), \texttt{alternative="two-sided"}, and
default \texttt{method="auto"}. With $n=1{,}576$ pairs
($n>50$), SciPy selects the \emph{asymptotic} normal
approximation; we do not apply a continuity correction
(\texttt{correction=False}). Reported $p$-values are therefore
asymptotic, not exact enumerations of the signed-rank null.
We report these $p$-values (in scientific notation or as
conservative powers-of-ten bounds) rather than only fixed-$\alpha$
accept/reject decisions.

\paragraph{Cohen's $d$ (paired).} The standardized effect size,
defined as the mean of paired differences divided by their
standard deviation (Bessel-corrected). Conventional
interpretation thresholds are $|d| \geq 0.8$ (large),
$0.5 \leq |d| < 0.8$ (medium), $0.2 \leq |d| < 0.5$ (small). We
report this alongside significance to disentangle effect magnitude
from sample-size-driven significance.

\subsection{Computing Environment}\label{sec:setup-environment}

Experiments were executed on a Python 3 environment with
scikit-learn $\geq$1.4, pandas $\geq$2.1, NumPy $\geq$1.26,
joblib $\geq$1.3 for parallelism, River $\geq$0.21 for ADWIN, and
SciPy $\geq$1.12 for statistical tests. Dependency floors are
documented in the accompanying artifact's
\texttt{requirements.txt}. Runs were parallelized across 8 worker
processes via joblib. Wall-clock times for the headline-producing
runs were: per-VM forecaster training across the 1{,}576-VM
corpus, 783 seconds; full-corpus profile evaluation (9{,}456
method-VM runs across six methods---two proposed and four
baselines), 1{,}162 seconds; corpus normalized-churn rerun for
Periodic-Predictor and DynAPM-V3, 297 seconds. The provisioning
evaluation was bounded by the joinable parquet replays and
completed within approximately 20 minutes wall-clock. The total
compute budget for the headline results was approximately one hour
on a single workstation with 8 worker processes.

\paragraph{AI-assisted research tooling.}
Consistent with the ACM Policy on Authorship, we disclose generative
AI use in the research workflow. Anthropic's Claude, accessed through
the Cursor development environment, assisted with implementing and
editing evaluation-pipeline code and with interpreting empirical
outputs during analysis. AI tools did not define the research
questions, experimental design, or scientific claims. All
AI-assisted code and analysis were reviewed and verified by the
authors against executable runs and project tests; the authors remain
solely responsible for the correctness of the reported results.

% ============================================================
% 5. Results
% ============================================================
\section{Results}\label{sec:results}

This section presents the empirical results of evaluating the two
proposed methods (Periodic-Predictor and DynAPM-V3) and four
observation-window baselines on 1{,}576 production VMs from the
Materna data-center corpus. Section~\ref{sec:results-accuracy} reports the
headline profile-accuracy results at corpus scale.
Section~\ref{sec:results-pairwise} presents pairwise statistical
tests for the two contributions, predictor-grounded refresh and
drift augmentation. Section~\ref{sec:results-tier} reports
per-volatility-tier analysis showing where each contribution is
most consequential. Section~\ref{sec:results-provisioning}
evaluates the downstream resource-provisioning task.
Section~\ref{sec:results-ablation} characterizes the drift
augmentation in more detail through the V3-vs-PP per-VM
distribution and reports normalized refresh-count statistics.

\subsection{Profile Accuracy at Corpus Scale}\label{sec:results-accuracy}

Table~\ref{tab:method-comparison} summarizes profile accuracy and
refresh activity for all six methods on the full 1{,}576-VM
corpus. We report mean CPU and memory class accuracy, mean p95
absolute error, and mean refresh-event count per VM. Each
metric is reported with a 95\% confidence interval computed via
1{,}000-sample paired bootstrap over the VM index.

\begin{table*}
\centering
\caption{Per-method profile accuracy and refresh activity on the
full Materna corpus (1{,}576 VMs; six methods). Mean values are
reported with 95\% confidence intervals from 1{,}000-sample paired
bootstrap. Rows are ordered alphabetically by method identifier
(DynAPM-V3, Full-Rewrite, Periodic-Window, Periodic-Predictor,
Reactive, Static), not by accuracy.}
\label{tab:method-comparison}
\input{tables/table1_method_comparison.tex}
\end{table*}

The headline observation is that the two methods using
predictor-grounded refresh (DynAPM-V3 and Periodic-Predictor) are
clearly separated from the four observation-window methods. CPU
class accuracy for DynAPM-V3 is 0.870 (95\% CI [0.863, 0.878]) and
for Periodic-Predictor is 0.868 (95\% CI [0.860, 0.876]). The best
observation-window method, Full-Rewrite, achieves 0.758
(95\% CI [0.748, 0.767]). The 11.2 percentage point gap between
Full-Rewrite and Periodic-Predictor is substantial despite
Full-Rewrite's much higher refresh activity (1{,}253 refresh
events per VM compared to Periodic-Predictor's 85.4): more
refreshes do not compensate for grounding refreshes in observed
history rather than in forecaster output.

The same pattern holds for memory class accuracy (DynAPM-V3 0.893,
Periodic-Predictor 0.890, Full-Rewrite 0.763) and for p95
prediction error (DynAPM-V3 1.04 percentage points CPU MAE,
Periodic-Predictor 1.06; Full-Rewrite 1.29; Periodic-Window 1.40).

The Static baseline, which never updates the profile after
deployment, provides the floor on what static profile
representations can achieve: 0.475 CPU class accuracy and 2.56
percentage points CPU p95 MAE. The Reactive baseline performs
worse than several scheduled-refresh baselines on the class
metric (0.645 CPU accuracy versus 0.737 for Periodic-Window),
illustrating that observation-window refresh combined with
threshold-triggered timing chatters under noise without yielding
a corresponding accuracy benefit.

Figure~\ref{fig:method-progression} visualizes the per-method CPU
class accuracy ranking. The visual gap between the four
observation-window methods (Static, Reactive, Periodic-Window,
Full-Rewrite) and the two predictor-grounded methods
(Periodic-Predictor, DynAPM-V3) is the central qualitative
result: predictor-grounded refresh is not a marginal optimization
but a categorically different operating point on the
profile-maintenance design space.

\begin{figure}
  \centering
  \includegraphics[width=\linewidth]{figs/figure1_method_progression.pdf}
  \caption{CPU class accuracy by method, sorted ascending.
Error bars represent 95\% bootstrap confidence intervals over per-VM
means ($n_\mathrm{vms}=1{,}576$). The two predictor-grounded methods
are visibly separated from the four observation-window baselines.}
  \Description{Vertical bar chart of mean CPU class accuracy for six
methods ordered left to right from lowest to highest accuracy
(Static, Reactive, Periodic-Window, Full-Rewrite, Periodic-Predictor,
DynAPM-V3), each with a small 95\% bootstrap CI error bar. Accuracy
rises across the four observation-window methods, then jumps to a
higher, nearly tied pair for Periodic-Predictor and DynAPM-V3.}
  \label{fig:method-progression}
\end{figure}

\subsection{Pairwise Statistical Tests}\label{sec:results-pairwise}

Table~\ref{tab:pairwise-deltas} reports paired statistical
comparisons for the two contributions. Each row reports the mean
per-VM difference between two methods, its 95\% bootstrap
confidence interval, the two-sided Wilcoxon signed-rank
$p$-value, and Cohen's $d$ for paired samples. All comparisons
use $n=1{,}576$ paired VMs.

\begin{table*}
\centering
\caption{Pairwise per-VM differences between methods. Mean
deltas, 95\% bootstrap CIs (1{,}000 resamples), two-sided
Wilcoxon signed-rank $p$-values from SciPy's asymptotic
approximation (\texttt{zero\_method="wilcox"}; zeros discarded;
no continuity correction), and Cohen's $d$ for paired samples.
All rows use $n=1{,}576$ VMs.}
\label{tab:pairwise-deltas}
\input{tables/table2_pairwise_deltas.tex}
\end{table*}

\paragraph{Predictor-grounded refresh.} The first four rows of the
table compare Periodic-Predictor against Periodic-Window. CPU
class accuracy improves by 13.05 percentage points
(95\% CI [12.08, 14.03], asymptotic Wilcoxon
$p = 5.17 \times 10^{-148}$, Cohen's $d = 0.69$). Memory class
accuracy improves by 14.54 percentage points
(95\% CI [13.71, 15.31], asymptotic Wilcoxon
$p = 1.50 \times 10^{-190}$, Cohen's $d = 0.89$). Both effects
are statistically robust at any conventional significance
threshold and have effect sizes in the medium-to-large range.
The CPU and memory p95 prediction errors also decrease (mean
reduction of 0.34 percentage points CPU,
$p = 1.20 \times 10^{-42}$; 0.52 percentage points memory,
$p = 2.38 \times 10^{-22}$), though their Cohen's $d$ values are
smaller (0.12 and 0.25), as expected for an absolute-error
metric whose per-VM variance is larger than for the bounded
class-accuracy metric.

The narrowness of the change distinguishes this from typical ML
methods comparisons. We do not introduce a new forecasting model
or change the refresh schedule. The single change is the
\emph{source} of the update content: forecaster prediction rather
than recent observation. The 13-to-14 percentage point gap is the
empirical magnitude of this change at corpus scale.

\paragraph{Drift augmentation.} The next four rows compare
DynAPM-V3 against Periodic-Predictor. CPU class accuracy improves by 0.24
percentage points (95\% CI [0.19, 0.28], asymptotic Wilcoxon
$p = 2.76 \times 10^{-18}$, Cohen's $d = 0.23$). Memory class accuracy improves by 0.32
percentage points ($d = 0.28$). p95 prediction errors decrease
by 0.02 percentage points each. The improvements are small in
absolute terms and small in effect size, but the bootstrap
confidence intervals exclude zero and the $p$-values are highly
significant. This is the expected pattern when a refinement is
genuine but limited in scope: the refinement helps most VMs by a
small amount, with the consistency of direction generating the
significance.

Figure~\ref{fig:predictor-lift} shows the per-VM distribution of
the predictor-grounded effect (CPU class accuracy of
Periodic-Predictor minus CPU class accuracy of Periodic-Window).
The mean improvement is $+0.130$ (median $0.125$), with the
distribution mass-shifted above zero; a minority of VMs are
negative, including a thin left tail of larger reversals, while
most paired differences remain positive.

\begin{figure}
  \centering
  \includegraphics[width=\linewidth]{figs/figure2_predictor_lift_distribution.pdf}
  \caption{Per-VM distribution of the predictor-grounded effect on
CPU class accuracy: $\mathrm{Periodic\text{-}Predictor} -
\mathrm{Periodic\text{-}Window}$. Each histogram bar represents
VM count. Vertical lines indicate the corpus mean and median.
The distribution is mass-shifted positive across the corpus.}
  \Description{Histogram of per-VM CPU class-accuracy differences
(Periodic-Predictor minus Periodic-Window) for 1{,}576 VMs. The
horizontal axis is the paired difference; the vertical axis is VM
count. Most mass lies above zero, with dashed mean and dotted median
reference lines both near $+0.13$. A thinner left tail of negative
differences extends below zero.}
  \label{fig:predictor-lift}
\end{figure}

\subsection{Per-Volatility-Tier Analysis}\label{sec:results-tier}

Table~\ref{tab:per-tier} reports CPU class accuracy and CPU
under-provisioning rate for each method, stratified by VM
volatility tier (low, medium, high). Tier sizes are 525, 525, and
526 VMs respectively (Section~\ref{sec:setup-volatility}). For
the predictor-grounded methods, the relative gain over
observation-window baselines grows with workload volatility.

\begin{table*}
\centering
\caption{Per-volatility-tier results. CPU class accuracy and CPU
under-provisioning rate for each method, by tier. The
predictor-grounded methods' relative gains are largest in the
high-volatility tier.}
\label{tab:per-tier}
\input{tables/table3_per_tier.tex}
\end{table*}

The CPU under-provisioning rate (computed as part of the
provisioning evaluation,
Section~\ref{sec:results-provisioning}) shows the most
operationally meaningful per-tier story. For the high-volatility
tier --- the 526 most variable VMs in the corpus, where
provisioning errors are most consequential ---
Periodic-Predictor cuts the under-rate from 14.79\% (under
Periodic-Window) to 9.08\%, an absolute reduction of 5.71
percentage points and a relative reduction of 38.6\%. For the
medium-volatility tier the absolute reduction is 1.71 percentage
points (relative 21.5\%). For the low-volatility tier
Periodic-Predictor's under-rate is marginally higher than
Periodic-Window's (3.07\% vs 2.83\%, an absolute increase of
0.24 percentage points), reflecting the fact that on stable
workloads, both methods are accurate enough that the residual
error is dominated by per-window noise rather than by
forecasting source.

This per-tier pattern is consistent with the operational
intuition behind the methodology: forecasting helps most where
recent observations are least reliable as a predictor of future
behavior. On low-volatility workloads, the previous 24 hours
predict the next 24 hours well, and observation-window refresh
suffices. On high-volatility workloads, the predictor's ability
to integrate longer-range features (lagged values, day-of-week
seasonality, rolling-window statistics over multiple horizons)
yields substantial accuracy gains that compound into significant
reductions in operational SLA-risk.

Figure~\ref{fig:per-tier-under-rate} visualizes the
under-provisioning rates for all six methods broken down by
tier. The visual clarity of the high-tier gap between
predictor-grounded and observation-window methods is the most
compelling figure for the operational impact of the
methodology.

\begin{figure*}
\centering
\includegraphics[width=\linewidth]{figs/figure4_per_tier_under_rate.pdf}
\caption{Mean CPU under-provisioning rate by method, grouped by
volatility tier. Predictor-grounded methods (Periodic-Predictor
and DynAPM-V3) cut the high-volatility-tier under-rate by 39\%
relative to Periodic-Window.}
\Description{Grouped bar chart of mean CPU under-provisioning rate.
Methods appear along the horizontal axis; within each method, three
bars show low-, medium-, and high-volatility tiers. Under-rates rise
with volatility for every method. Reactive is highest on the
high-volatility tier; Periodic-Predictor and DynAPM-V3 are nearly
tied and lowest on that tier, below Periodic-Window and Full-Rewrite.}
\label{fig:per-tier-under-rate}
\end{figure*}

\subsection{Resource Provisioning Recommendation Task}\label{sec:results-provisioning}

The provisioning task (Section~\ref{sec:setup-provisioning})
treats each profile snapshot as input to a downstream LCM
consumer that recommends a provisioned capacity of $1.2 \times
\texttt{predicted\_cpu\_p95\_next\_24h}$. We measure the rate of
under-provisioning (actual p95 exceeds recommended capacity, an
SLA-risk event) and over-provisioning (recommended capacity
exceeds 1.5 times actual p95, a wasted-capacity event) at each
evaluation timestamp.

Figure~\ref{fig:provisioning-outcomes} shows stacked
under/well/over rates for all six methods. Predictor-grounded
methods achieve the highest well-provisioned share (DynAPM-V3
88.2\%, Periodic-Predictor 87.8\%); the remaining 12\% is split
roughly evenly between under-provisioning (6.1\% for both
methods) and over-provisioning (5.7\% V3, 6.0\% PP). Static and
Reactive sit at the operationally bad extremes: Static
over-provisions in 16.0\% of windows (worst over-rate) while
Reactive under-provisions in 16.4\% (worst under-rate).

\begin{figure}
  \centering
  \includegraphics[width=\linewidth]{figs/figure3_provisioning_outcomes.pdf}
  \caption{Resource provisioning outcomes for all six evaluated methods
at safety factor 1.2. Stacked bars show the fraction of evaluation
points classified as under-provisioned (SLA risk), well-provisioned,
or over-provisioned (wasted capacity), averaged within each VM and
then across VMs.}
  \Description{Stacked bar chart for six methods at safety factor 1.2.
Each bar sums to one and stacks under-provisioned (bottom),
well-provisioned (middle), and over-provisioned (top) fractions after
averaging within then across VMs. Periodic-Predictor and DynAPM-V3
show the largest well-provisioned middle segments and small under
and over caps; Static has the largest over segment; Reactive has the
largest under segment.}
  \label{fig:provisioning-outcomes}
\end{figure}

\paragraph{Predictor-grounded versus observation-window provisioning.}
Comparing Periodic-Predictor against Periodic-Window: CPU
under-rate drops from 8.52\% to 6.13\%, an absolute reduction of
2.39 percentage points and a relative reduction of 28\%
(95\% bootstrap CI on the absolute delta:
[$-3.07\%, -1.67\%$], asymptotic Wilcoxon
$p = 7.06 \times 10^{-62}$, Cohen's
$d = -0.17$). Memory under-rate similarly drops by 1.55
percentage points ($p = 1.86 \times 10^{-57}$). The over-rate change is
small for CPU but substantial for memory ($-3.20$ percentage
points absolute, $p = 4.77 \times 10^{-83}$), reflecting the larger
forecasting improvement on memory than on CPU.

\paragraph{Drift augmentation in provisioning.} DynAPM-V3 ties
Periodic-Predictor on under-rate (the difference is 0.016
percentage points, with a CI containing zero and Wilcoxon
$p = 0.21$, statistically not distinguishable from no effect),
but reduces over-rate by 0.31 percentage points
($p = 1.24 \times 10^{-26}$, Cohen's $d = -0.30$). The interpretation is that drift
augmentation does not measurably reduce SLA-risk events further
beyond what predictor-grounded refresh already achieves, but it
does reduce wasted capacity by trimming over-allocation when the
forecaster's residuals indicate the workload has shifted to a
lower regime.

\paragraph{Sensitivity to safety factor.} We re-evaluated the
provisioning task at safety factors $s \in \{1.1, 1.2, 1.5\}$ to
verify that the method ranking is not an artifact of the chosen
safety margin. Across all three values of $s$, the ordering of
methods by under-rate is preserved: predictor-grounded methods
are always the lowest, observation-window scheduled methods
intermediate, and Reactive is always worst. The
predictor-grounded under-rate advantage over Periodic-Window has
absolute magnitudes of 1.4, 2.4, and 1.8 percentage points at
$s = 1.1, 1.2, 1.5$ respectively. The qualitative ranking of
methods is therefore robust to the choice of safety margin in
the range we evaluated.

\subsection{Drift Augmentation: Per-VM Distribution and Refresh Activity}\label{sec:results-ablation}

To characterize when drift augmentation helps, we examine the
per-VM distribution of the V3-versus-PP CPU class accuracy
delta. Figure~\ref{fig:v3-vs-pp-scatter} plots each VM as a point
with x-coordinate equal to Periodic-Predictor's CPU class
accuracy and y-coordinate equal to DynAPM-V3's accuracy on the
same VM. Points above the $y=x$ diagonal are VMs where drift
augmentation helps; points below are VMs where it hurts.

\begin{figure}
  \centering
  \includegraphics[width=\linewidth]{figs/figure5_v3_vs_pp_per_vm_corpus.pdf}
  \caption{Per-VM CPU class accuracy: DynAPM-V3 versus
Periodic-Predictor (n=1{,}576). Points above the y=x diagonal
indicate VMs where drift augmentation improves accuracy; points
below indicate VMs where it hurts. The corpus distribution shows
a strong concentration on the diagonal (drift augmentation
neutral on most VMs), with a long upper tail of VMs that benefit
substantially from the additional refreshes.}
  \Description{Scatter plot of 1{,}576 VMs with Periodic-Predictor CPU
class accuracy on the x-axis and DynAPM-V3 accuracy on the y-axis,
plus a dashed $y=x$ line. Most points lie on the diagonal (equal
accuracy). Points colored as V3 greater than PP sit above the line,
including a sparse upper tail of larger gains; only a few points fall
slightly below the line.}
  \label{fig:v3-vs-pp-scatter}
\end{figure}

The distribution is highly asymmetric. Of the 1{,}576 VMs:
\begin{itemize}
\item 1{,}427 VMs (90.5\%) have V3 = PP exactly --- drift
augmentation triggers no improvement, typically because no drift
events fire during the operational period or because the periodic
schedule already captures the relevant changes.
\item 135 VMs (8.6\%) have V3 \textgreater\ PP --- drift events fired
and the additional refreshes improved accuracy.
\item 14 VMs (0.9\%) have V3 \textless\ PP --- drift events fired
and the additional refreshes degraded accuracy.
\end{itemize}

The ratio of helped to hurt VMs is approximately 10:1, and the
maximum per-VM gain (12.5 percentage points) is more than twice
the maximum per-VM loss (5.3 percentage points). The per-VM mean
gain of 0.24 percentage points is therefore not characteristic of
the corpus as a whole: it is concentrated in a minority of VMs,
mostly in the high-volatility tier, where drift augmentation
delivers operationally meaningful improvements.

\paragraph{Refresh activity.} DynAPM-V3 commits 91.3 refresh events
per VM on average across the operational period (median 89),
compared to Periodic-Predictor's 85.4 (median 86). The
decomposition for V3 is 85.4 periodic refreshes plus 5.96 drift
refreshes per VM on average. The paired per-VM delta is 5.96
events with 95\% bootstrap CI [5.57, 6.33], asymptotic Wilcoxon
$p = 1.54 \times 10^{-212}$. The drift augmentation increases refresh activity by
about 7\% on average; this overhead is the cost paid for the
gains documented in this section. Methods that refresh more
aggressively (Full-Rewrite at 1{,}253 events per VM) achieve
substantially worse accuracy, confirming that refresh frequency
is not the primary lever and that the source of update content
matters more than the timing of updates.

% ============================================================
% 6. Discussion
% ============================================================
\section{Discussion}\label{sec:discussion}

This section discusses the implications of the two contributions,
identifies threats to validity, and considers the conditions
under which the findings reported in Section~\ref{sec:results}
would and would not generalize beyond our experimental setting.

\subsection{Implications of Predictor-Grounded Refresh}\label{sec:discussion-pp}

The empirical magnitude of the predictor-grounded refresh
effect --- 13.05 percentage points CPU class accuracy and 14.54
percentage points memory class accuracy over a matched-schedule
observation-window baseline --- is not a marginal improvement.
At this magnitude on a 1{,}576-VM corpus, with bootstrap
confidence intervals on the order of one percentage point and
$p$-values below $10^{-147}$, the effect is qualitatively distinct
from the kind of accuracy lift achievable through tuning of an
observation-window method. The 11.2 percentage point gap between
Periodic-Predictor and Full-Rewrite, despite Full-Rewrite's
1{,}253 refreshes per VM versus Periodic-Predictor's 85, indicates
that this gap is not closable by increasing refresh frequency
within the observation-window paradigm.

The mechanism explaining this gap is straightforward in retrospect.
Profile dynamic fields are read by downstream consumers that act
on \emph{future} behavior: an autoscaler provisions for the
upcoming window, a placement scheduler considers projected load,
a capacity planner reasons about the next operational period.
When dynamic fields summarize past observations, the consumer must
either accept that summary as a proxy for the future (introducing
prediction error implicit in the proxy assumption) or perform its
own forecasting on top of the profile (fragmenting forecasting
logic across multiple LCM components). Predictor-grounded refresh
moves the forecasting step into the profile-maintenance layer
itself: the per-VM forecaster is consulted at each refresh, and
its output becomes the profile's dynamic-field content. Consumers
read forecasts directly. The 13--15 percentage point accuracy gap
quantifies the information loss imposed by the
observation-as-proxy assumption.

This finding has design implications for application-profile
schemas intended for CECC LCM systems. Schemas that explicitly
distinguish between observed and predicted dynamic fields, and
that provide first-class support for forecast-grounded updates,
should produce more actionable profiles than schemas that treat
dynamic fields as monitoring summaries---at least for
data-center/cloud-tier consumers of the kind evaluated here.
Concretely, schema designs should consider semantic versioning of
dynamic fields (observed-window vs.\ predicted-window), explicit
forecast-horizon metadata, and update-event provenance so that
consumers know which version of a field they are reading.

\subsection{Why Drift Augmentation Provides Small but Consistent Gains}\label{sec:discussion-drift}

The drift-augmentation contribution is genuine but smaller in
scope than the predictor-grounded contribution. The corpus-mean
gain is 0.24 percentage points CPU class accuracy with $p$ below
$10^{-17}$, but the per-VM distribution
(Section~\ref{sec:results-ablation}) reveals that 90.5\% of VMs
see no change at all from drift augmentation, while 8.6\% of VMs
benefit, and 0.9\% are slightly hurt. The corpus mean is a
weighted average across these three populations.

Two observations explain this pattern. First, the periodic
refresh schedule already captures most of the actionable change
in workload behavior: at $\Delta_p = 6$ hours, scheduled refreshes
occur 4 times per day and four times more frequently than the
forecast horizon's 24-hour scope. Drift events between scheduled
refreshes correspond to unexpectedly fast change relative to this
already-frequent baseline. Such events are sparse on most VMs
because most VMs in production data centers are stable enough
that 6-hour periodic refresh suffices.

Second, when drift events do fire, they correspond to
operationally meaningful change. The 135 VMs where DynAPM-V3
improves over Periodic-Predictor are concentrated in the
high-volatility tier (Section~\ref{sec:results-tier}); the
maximum per-VM gain is 12.5 percentage points, more than twice
the maximum loss of 5.3 percentage points. The asymmetry is
consistent with the operational interpretation of the mechanism:
ADWIN signals drift only when the residual stream contains
statistically distinct sub-windows, and the ensuing refresh
incorporates the most recent signal into the profile. On VMs
where the workload is genuinely changing, this targeted
intervention helps. On VMs where the workload is stable, ADWIN
seldom signals, and DynAPM-V3 reduces to Periodic-Predictor.

The result also tempers an instinctive expectation that adding
``intelligence'' to a maintenance method should yield large
gains. In our setting, the largest gain came from rethinking the
\emph{source} of the update (predictor versus observation),
not from elaborating the trigger logic with a sophisticated
detector. The methodological lesson is that data-source decisions
should be exhausted before trigger-mechanism elaboration is
considered.

\subsection{Threats to Validity}\label{sec:discussion-threats}

We identify four threats to validity for the findings reported
in Section~\ref{sec:results}.

\paragraph{External validity (corpus and deployment tier).}
Our evaluation uses the Materna VM trace exclusively, under
offline replay of CPU and memory utilization. Materna is a
substantial corpus (1{,}576 VMs across three monthly snapshots of
heterogeneous business applications) but represents a single
deployment context: a European enterprise data center on VMware
ESX---i.e., the data-center/cloud tier of the continuum, not a
full CECC testbed. We do not evaluate edge-tier workloads,
containerized or Kubernetes deployments, heterogeneous
cloud--edge platforms, or live orchestrator integration.
Public-cloud traces (e.g., Google Borg or Alibaba) and edge-tier
workloads may exhibit different volatility, periodicity, and
scaling behavior. The qualitative finding that predictor-grounded
refresh outperforms observation-window refresh should transfer
wherever workloads are non-trivially predictable from
per-instance history, but the magnitude of the effect may differ,
and continuum-wide claims require additional validation.

\paragraph{Single forecaster.} We use RandomForest with fixed
hyperparameters as the per-VM forecaster. The forecaster is not
the contribution of this work, but its choice affects the
absolute accuracy of predictor-grounded methods. A weaker
forecaster would likely shrink the predictor-grounded gap; a
stronger one would likely widen it. We did not compare alternative
forecasters in this study, so we do not claim robustness of the
gap magnitude to model family. The qualitative finding that
predictor-sourced updates outperform matched-schedule
observation-window updates presupposes that the forecaster
carries some signal about future behavior on the evaluated
traces.

\paragraph{Single class scheme.} The five-class label space
(\textsc{idle}, \textsc{low}, \textsc{medium}, \textsc{high},
\textsc{burst}) and the per-VM tertile-based volatility tiers
are specific to our schema. The accuracy advantage of
predictor-grounded refresh applies to any classification scheme
where the predicted distribution carries information beyond the
observed distribution; the specific percentage-point magnitudes
are sensitive to bin boundaries (a four-class scheme with wider
bins would naturally have higher absolute accuracy and a smaller
percentage-point gap, even though the underlying mechanism is
unchanged).

\paragraph{Snapshot-and-verify protocol versus production deployment.}
Our evaluation protocol reads profile snapshots at fixed
12-hour grid points and verifies them against ground truth from
the subsequent 24 hours. This protocol abstracts away the
operational behaviors of real LCM consumers, which may read
profiles at unpredictable times, may persist their own copies
of profiles, and may interact with profile maintenance
asynchronously. The protocol is appropriate for measuring profile
representation accuracy as an isolated property; integrating
DynAPM into a production LCM stack and measuring end-to-end
operational metrics (e.g., actual SLA-violation rates, resource
utilization, cost) is future work.

\subsection{Generalization}\label{sec:discussion-generalization}

The methodology is intended for CECC LCM systems that share two
structural features: applications with per-instance behavioral
heterogeneity (so that per-application forecasters carry
meaningful signal), and consumers that read profile snapshots and
act on near-future behavior. Our evidence for those features comes
from production data-center VMs; we hypothesize---but do not
demonstrate---that the same design applies at edge and hybrid
tiers when analogous telemetry and profile consumers exist.

The contributions are likely \emph{not} useful where applications
are too homogeneous for per-application forecasting to beat global
heuristics, where behavior is effectively unpredictable, or where
downstream consumers do not read the profile fields whose accuracy
we measure. The third case is a deployment concern that warrants
integration testing before any production rollout.

Replication on public-cloud, edge-tier, containerized, and live
orchestrated settings is needed to validate effect magnitudes
beyond the Materna data-center tier.

% ============================================================
% 7. Conclusion and Future Work
% ============================================================
\section{Conclusion and Future Work}\label{sec:conclusion}

We presented DynAPM, a profile-maintenance methodology intended for
lifecycle management in the cloud-edge continuum, in which dynamic
profile fields encode forecasts of near-future application behavior
rather than summaries of recent observations. Within this framework
we propose two methods: Periodic-Predictor, which performs
predictor-grounded refresh on a fixed schedule, and DynAPM-V3, which
augments that schedule with drift-triggered refreshes on the
prediction residual stream. Evaluating both proposed methods against
four observation-window baselines (Static, Periodic-Window, Reactive,
and Full-Rewrite)---six methods in total---via offline replay on
1{,}576 production VMs from the Materna enterprise data-center
corpus (the data-center/cloud tier), we found that predictor-grounded
refresh substantially improves profile accuracy over a
matched-schedule observation-window baseline, that this accuracy gain
translates into a corresponding reduction in SLA-risk events on a
downstream provisioning task, and that drift augmentation yields a
small but statistically robust further refinement concentrated on
high-volatility workloads.

The methodological lesson from this work is narrower and more
specific than ``prediction beats observation.'' The lesson is
that the \emph{source} of profile updates --- whether observed
history or model-based forecast --- is a first-order design
choice for profile-maintenance methods, separable from the
trigger logic that decides when to refresh. Both academic and
production profile-maintenance designs have largely treated the
update source as an implicit by-product of the trigger
mechanism. Treating the source as a primary design dimension,
and grounding it in per-application forecasters, opens a
previously underexplored region of the design space.

Several directions follow naturally from the limitations
identified in Section~\ref{sec:discussion-threats}.
Cross-dataset replication on public-cloud traces and, separately,
edge-tier, containerized, and heterogeneous cloud--edge workloads
would test whether the predictor-grounded effect holds beyond the
Materna data-center tier. End-to-end integration with a live LCM
or orchestrator stack, with measurement of operational SLA,
utilization, and cost outcomes, remains future work. Substituting
more expressive forecasters---attention-based architectures,
ensemble models, or online-learning forecasters that adapt to
drift directly---would characterize the upper bound on what
predictor-grounded refresh can achieve. Finally, the asymmetric
per-VM distribution of the drift-augmentation effect
(Section~\ref{sec:results-ablation}) suggests that selective
drift-augmentation policies---enabling drift refreshes only for
VMs whose volatility characteristics indicate the mechanism will
help---could improve the cost-benefit profile of DynAPM-V3 further.

% ============================================================
% Declarations
% ============================================================

\section*{Declaration of Competing Interest}
The authors declare that they have no known competing financial interests
or personal relationships that could have appeared to influence the work
reported in this paper.

\section*{Declaration of Generative AI and AI-Assisted Technologies}
We used Anthropic's Claude through the Cursor environment as follows.
\emph{Code assistance:} drafting and editing evaluation-pipeline and
supporting scripts. \emph{Analysis assistance:} help interpreting
empirical outputs and organizing comparisons. \emph{Language and
drafting assistance:} help editing and drafting manuscript prose
(disclosed for transparency; ACM does not currently require
disclosure of writing assistance alone). AI systems are not authors.
The authors reviewed and revised all AI-assisted code, analyses, and
text, verified empirical claims against executed experiments and
tests, and take full responsibility for the content of this article.
Research uses of AI are also summarized in
Section~\ref{sec:setup-environment}.

\section*{Data Availability}
The Materna VM utilization traces used in this study are publicly
available from the Grid Workloads Archive~\cite{iosup_gwa_2008,materna_trace}
and via Kaggle at
\url{https://www.kaggle.com/datasets/kpiyush04/maternaworkloadtraces}.
Raw Materna CSVs are not redistributed with the submission artifact.
Evaluation code, scripts, and regenerated paper tables/figures will be
made available via an anonymous artifact link for review
(\texttt{ANONYMOUS-ARTIFACT-URL}) and via a public archival release upon
acceptance.

\section*{Artifact Availability}
An anonymous review artifact will be hosted at
\texttt{ANONYMOUS-ARTIFACT-URL}. The intended contents include source
code and scripts to regenerate the reported tables and figures from
precomputed intermediate artifacts, dependency pins
(\texttt{requirements.txt}), and documentation sufficient to rerun the
evaluation pipeline. Large raw traces and full intermediate result
trees may be omitted for size; instructions will describe how to obtain
Materna data from the Grid Workloads Archive and which derived
artifacts are included.

% ============================================================
% Bibliography
% ============================================================
\bibliographystyle{ACM-Reference-Format}
\bibliography{cas-refs}

\end{document}